\bookchapter{الفصل 18: النسخ الاحتياطي للبيانات وترحيلها}

\chapterintro{جارتي، سقط هاتفها في حوض الماء وأخرجته ولم يعد يعمل. لم يكن ألمها للهاتف نفسه، بل للصور التي لم تنسخها احتياطياً منذ ثلاث سنوات — كل سجل طفلها من الولادة حتى المشي. بعد سماع قصتها شعرت بسوء شديد، لأن هذا كان يمكن تجنبه تماماً. في هذا الفصل سنؤمّن بياناتك بشكل جدي.}

\booksection{أولاً: اعرف أولاً: ما الذي إذا فُقد من هاتفك لا يعود}

“حالة البقاء” للبيانات المختلفة ليست متشابهة، يجب أن تميز بينها:

\begin{itemize}
  \item \textbf{مخزن محلياً على الهاتف، وإذا فُقد فقد زال}: الصور والفيديوهات (إذا لم تُفعّل المزامنة السحابية)، والملاحظات، والرسائل النصية، وسجل المكالمات، والملفات المنزّلة.
  \item \textbf{مخزن على خوادم التطبيقات، وعند تغيير الهاتف وإعادة تسجيل الدخول تجده}: قائمة أصدقاء WeChat، وطلبات Taobao، وسجل المعاملات في تطبيقات البنوك. هذه الأشياء مخزنة على خوادم الشركات المعنية، وليست على هاتفك.
  \item \textbf{في منطقة رمادية، تتطلب إجراءات خاصة}: سجل محادثات WeChat (رغم أن خادم WeChat يخزنه، لكنك لا تستطيع تنزيله من الخادم)، وبيانات اعتماد الأجهزة المرتبطة بتطبيقات البنوك، والمفاتيح في تطبيقات المصادقة الثنائية (مثل Google Authenticator).
\end{itemize}

لذا، ما يجب أن تحرص على حمايته هو الفئة الأولى والثالثة.

\booksection{ثانياً: النسخ الاحتياطي الكامل للجهاز: احصل على كل شيء دفعة واحدة}

النسخ الاحتياطي الكامل للجهاز مثل التقاط “لقطة” كاملة للهاتف — إعدادات النظام الحالية، وتخطيط سطح المكتب، وبيانات التطبيقات، والصور، وجهات الاتصال، كلها معاً.

\booksubsection{الخيار الأول: النسخ الاحتياطي المحلي على الحاسوب (الأكثر موثوقية، يُوصى به لمن لديهم حاسوب)}

\begin{itemize}
  \item \textbf{iPhone}: أوصل الهاتف بالحاسوب عبر كابل بيانات، افتح iTunes (الإصدار القديم لـ Windows) أو Finder (Mac)، وفي صفحة الجهاز اختر "النسخ الاحتياطي الآن". انتبه لتفعيل "تشفير النسخة الاحتياطية"، وإلا لن تُنسخ كلمات المرور والبيانات الصحية.
  \item \textbf{أندرويد}: لكل علامة برنامج خاص على الحاسوب، مثل Huawei Phone Assistant وXiaomi Assistant وSamsung Smart Switch. أوصل بالحاسوب واتبع التعليمات.
  \item ميزة النسخ الاحتياطي على الحاسوب أنه لا يتطلب شراء مساحة سحابية، وسرعة الاستعادة أسرع بكثير من الاستعادة السحابية. العيب هو أنه يتطلب منك التذكر الدوري لتوصيل الكابل والنسخ — ومعظم الناس لا يفعلون ذلك.
\end{itemize}

\booksubsection{الخيار الثاني: النسخ الاحتياطي السحابي (الأكثر راحة، يُوصى به للجميع)}

\begin{itemize}
  \item \textbf{iPhone}: الإعدادات اسمك iCloud iCloud Backup فعّل. عندما يكون الهاتف موصولاً بالطاقة وعلى شبكة Wi-Fi ومقفلاً، سيتم النسخ الاحتياطي تلقائياً. الـ 5 غيغابايت المجانية غالباً لا تكفي، يُنصح بدفع 6 يوانات للترقية إلى 50 غيغابايت.
  \item \textbf{أندرويد}: لكل علامة سحابتها ميزة النسخ الاحتياطي. اذهب لهواتف Huawei وابحث عن "النسخ الاحتياطي السحابي" في الإعدادات، Xiaomi إلى "Xiaomi Cloud Services"، OPPO/vivo مشابه. وهناك أيضاً نسخ احتياطي من Google مدمج (الإعدادات النظام النسخ الاحتياطي)، فعّله أيضاً.
  \item أكبر ميزة للنسخ الاحتياطي السحابي أنه "لا يتطلب منك القلق"، يتم تلقائياً. أكبر عيب هو أن الاستعادة تتطلب وقتاً طويلاً للتنزيل، والمساحة المجانية عموماً لا تكفي.
\end{itemize}

\booksection{ثالثاً: استراتيجية النسخ الاحتياطي المصنّف: الأشياء المهمة احفظ نسخة منفصلة}

النسخ الاحتياطي الكامل للجهاز هو الضمان الأساسي، لكن بعض الأشياء مهمة جداً وتستحق نسخة منفصلة:

\booksubsection{الصور والفيديوهات (الأصول الرقمية الأثمن)}

\begin{itemize}
  \item \textbf{خطة التأمين المزدوج}: فعّل المزامنة التلقائية في صور iCloud/Google Photos، وكل بضعة أشهر انسخ الصور يدوياً إلى الحاسوب أو قرص صلب خارجي. هكذا نسخة سحابية ونسخة محلية، كل منهما نسخة احتياطية للأخرى.
  \item \textbf{خطة بدون دفع}: Google Photos يمنحك مجاناً 15 غيغابايت، يمكن ضبطه على "جودة مضغوطة" لتوفير المساحة. أو كل شهر استخدم كابل بيانات لنسخ جميع الصور إلى الحاسوب.
\end{itemize}

\booksubsection{جهات الاتصال (يُسهى عنها لكنها بالغة الأهمية)}

\begin{itemize}
  \item جهات اتصال iPhone تتم مزامنتها افتراضياً مع iCloud، إذا فعّلت المزامنة فلا تحتاج لنسخ احتياطي إضافي. لكن إذا أردت مزيداً من الأمان، يمكنك تصدير ملف vCard من iCloud.com وحفظه على الحاسوب.
  \item جهات اتصال أندرويد تتم مزامنتها عادةً مع Google أو سحابة العلامة التجارية. اذهب إلى إعدادات تطبيق جهات الاتصال وصدّر ملف .vcf، فيكون لديك نسخة احتياطية كاملة.
\end{itemize}

\booksubsection{سجل محادثات WeChat (عالجه بشكل منفصل، الأهمية القصوى)}

يجب التحدث عن هذا بشكل منفصل، لأن النسخ الاحتياطي الكامل للجهاز كثيراً ما لا يشمل سجل محادثات WeChat بالكامل، وأدوات نقل الأجهزة لا تستطيع نقله أصلاً.

\begin{itemize}
  \item \textbf{الطريقة الأولى: النقل إلى هاتف جديد}. WeChat أنا الإعدادات عام ترحيل ونسخ سجل المحادثات احتياطياً ترحيل ترحيل إلى هاتف/جهاز لوحي. أوصل الهاتفين القديم والجديد بنفس شبكة Wi-Fi، وامسح الكود للنقل. انتبه: قبل اكتمال النقل، لا تخرج من WeChat على الهاتف القديم.
  \item \textbf{الطريقة الثانية: النسخ الاحتياطي إلى الحاسوب}. ثبّت WeChat PC على الحاسوب، الأشرطة الثلاثة في الأسفل يساراً الترحيل والنسخ الاحتياطي النسخ الاحتياطي والاستعادة النسخ الاحتياطي لسجل المحادثات إلى الحاسوب. لاحقاً عند تغيير الهاتف يمكنك الاستعادة إلى الهاتف الجديد.
  \item \textbf{الطريقة الثالثة: النسخ الاحتياطي إلى السحابة}. يدعم WeChat أيضاً النسخ الاحتياطي إلى iCloud أو المساحة السحابية للهاتف، لكن انتبه: هذا النسخ الاحتياطي يمكن استعادته فقط على هاتفك الحالي، ولا يمكن نقله عبر الأجهزة. لتغيير الهاتف، لا بد من استخدام الطريقتين السابقتين.
\end{itemize}

\booksection{رابعاً: دليل تغيير الهاتف: من الهاتف القديم إلى الجديد}

تغيير هاتف جديد أمر سار، لكن إذا لم تُرحّل البيانات فسيصبح صداعاً. ثلاث حالات تُعالج بشكل مختلف:

\booksubsection{الترقية ضمن العلامة نفسها (الأبسط)}

\begin{itemize}
  \item \textbf{iPhone iPhone جديد}: عند التشغيل والتفعيل، ضع الهاتف القديم بجانب الجديد، وستنبثق "البداية السريعة" تلقائياً، امسح النمط على الهاتف الجديد بالهاتف القديم. جميع البيانات والإعدادات وتخطيط التطبيقات تُرحّل كما هي، بما في ذلك كلمات مرور Wi-Fi. يستغرق حوالي 20-30 دقيقة.
  \item \textbf{أندرويد ضمن العلامة نفسها}: استخدم أداة نقل الأجهزة المدمجة من العلامة. Huawei تستخدم "Phone Clone"، Xiaomi تستخدم "Xiaomi Mover"، OPPO تستخدم "Phone Clone"، vivo تستخدم "EasyShare". الخطوات متشابهة: ثبّت التطبيق نفسه على الهاتفين القديم والجديد، امسح الكود للاتصال، حدد ما تريد نقله، ابدأ النقل.
\end{itemize}

\booksubsection{تغيير الهاتف عبر العلامات (أندرويد إلى أندرويد، علامات مختلفة)}

استخدم أداة نقل طرف ثالث أو أداة النقل المدمجة من العلامة القديمة (معظم أدوات نقل العلامات تدعم النقل إلى هواتف أندرويد أخرى). ما يمكن نقله: جهات الاتصال، الصور، الرسائل النصية، سجل المكالمات، بعض التطبيقات. ما لا يمكن نقله: حالة تسجيل الدخول داخل التطبيقات (عليك إعادة تسجيل الدخول)، سجل محادثات WeChat (استخدم ميزة الترحيل الخاصة بـ WeChat).

\booksubsection{تغيير الهاتف عبر المنصات (أندرويد إلى iPhone أو العكس)}

\begin{itemize}
  \item \textbf{أندرويد iPhone}: لدى Apple تطبيق "Move to iOS"، نزّله على هاتف أندرويد، وعند تفعيل iPhone الجديد اختر "نقل البيانات من جهاز أندرويد". يمكنه نقل جهات الاتصال والصور والرسائل النصية وغيرها من البيانات الأساسية، لكن بيانات التطبيقات لا يمكن نقلها. سجل محادثات WeChat لا بد من استخدام ميزة الترحيل الخاصة بـ WeChat.
  \item \textbf{iPhone أندرويد}: لكل علامة أندرويد أدوات نقل خاصة، لكن لنعترف، تجربة النقل عبر المنصات ليست جيدة كالتجربة ضمن المنصة نفسها. الأهم أن تنجز أولاً: تصدير جهات الاتصال، النسخ الاحتياطي للصور، ترحيل سجل محادثات WeChat. الباقي يمكن إنجازه تدريجياً.
\end{itemize}

\booksection{خامساً: القاعدة الذهبية للنسخ الاحتياطي: مبدأ الثلاثة-اثنين-واحد}

هذا مبدأ كلاسيكي في مجال أمن البيانات، بسيط وسهل التذكر:

\begin{itemize}
  \item \textbf{3}: احتفظ بثلاث نسخ على الأقل من البيانات المهمة (نسخة أصلية، نسختان احتياطيتان).
  \item \textbf{2}: استخدم وسائط تخزين مختلفة للنسخ الاحتياطي (مثلاً نسخة سحابية، نسخة على قرص صلب خارجي).
  \item \textbf{1}: ضع نسخة احتياطية واحدة على الأقل في موقع مادي مختلف (للحماية من الحرائق والسرقة التي قد تدمر جميع البيانات دفعة واحدة).
\end{itemize}

للشخص العادي، تحقيق “نسخة سحابية، نسخة على الحاسوب” يُعتبر جيداً جداً. لا تعقّد الأمور، المهم أن تنفّذها فعلاً.

\begin{tcbwarn}{انتبه بشكل خاص لسجل محادثات WeChat}
لا يمكن ترحيل سجل محادثات WeChat بأي أداة نقل هواتف! وهو ليس ضمن تغطية النسخ الاحتياطي الكامل للجهاز. يجب استخدام ميزة “ترحيل سجل المحادثات” المدمجة في WeChat، أو النسخ الاحتياطي إلى WeChat على الحاسوب. كثيرون يكتشفون بعد تغيير الهاتف الجديد أن سجل المحادثات قد ضاع كله، وعندئذ يكون الوقت قد فات — ربما يكون الهاتف القديم قد أُعيد ضبطه على المصنع.

\end{tcbwarn}

\begin{tcbpractice}{تدرّب: أجرِ تمرين نسخ احتياطي كامل}
\begin{enumerate}
  \item تحقق من تفعيل iCloud Backup أو النسخ الاحتياطي السحابي للعلامة التجارية، وتأكد من وقت آخر نسخة احتياطية.
  \item انسخ الصور يدوياً: أوصل الحاسوب بكابل بيانات، وانسخ مجلدي DCIM وPictures إلى الحاسوب.
  \item انسخ جهات الاتصال احتياطياً: مستخدمو iPhone يصدّرون vCard من iCloud.com؛ مستخدمو أندرويد يصدّرون ملف .vcf من إعدادات تطبيق جهات الاتصال.
  \item انسخ سجل محادثات WeChat احتياطياً: افتح WeChat أنا الإعدادات عام ترحيل ونسخ سجل المحادثات احتياطياً النسخ الاحتياطي إلى الحاسوب.
  \item انسخ ملف vCard لجهات الاتصال ومجلد الصور إلى قرص صلب خارجي أو فلاش USB، لتحقيق "نسخة واحدة غير متصلة في مكان آخر" ضمن مبدأ الثلاثة-اثنين-واحد.
\end{enumerate}

\end{tcbpractice}
